\section*{7 Orbit dimension, tangent spaces, and flows}
In this section we introduce tangent spaces and study their dimensions. We especially study them in the context of $G\fcom$, and $Gf_0$.
The orbit dimension of a tensor $f\in\C^{w_1w_2}\otimes\C^{w_2w_3}\otimes\cdots\otimes\C^{w_dw_{d+1}}$ is the dimension of the orbit $Gf$ as an affine variety. It can be determined as the dimension of the tangent space $T_f$ of the action of $G$ at $f$, which is a vector space defined as follows. Let $\g:=\C^{w_1w_2\times w_1w_2}\times\cdots\times\C^{w_dw_{d+1}\times w_dw_{d+1}}$. For $A\in\g$ we define the Lie algebra action $Af:=\lim_{\varepsilon\to0}\frac1\varepsilon((\mathrm{id}+\varepsilon A)f-f)$, where $\mathrm{id}\in G$ is the identity element. We define the vector space $T_f:=\g f=\{Af\mid A\in\g\}$.
\begin{claim}The dimension $\dim T_h$ is the same for all $h\in Gf$.\end{claim}
\begin{proof}
Since the action of $G$ is linear, for all $g\in G$ and $A\in\g$ we have
\begin{align*}
A(gf)&=\lim_{\varepsilon\to0}\frac1\varepsilon((\mathrm{id}+\varepsilon A)(gf)-gf)\\
&=\lim_{\varepsilon\to0}\frac1\varepsilon(gg^{-1}(\mathrm{id}+\varepsilon A)gf-gf)\\
&=g\lim_{\varepsilon\to0}\frac1\varepsilon((\mathrm{id}+\varepsilon(g^{-1}Ag))f-f)\\
&=g((g^{-1}Ag)f).
\end{align*}
Since $A\mapsto g^{-1}Ag$ is a bijection on $\g$, it follows that $T_{gf}=gT_f$. Hence the claim follows.
\end{proof}
In the following we will use Claim 10 to argue $f_0\notin G\fcom$ by showing that $\dim T_{\fcom}$ and $\dim T_{f_0}$ are different.
Let $e,e'\in E^i$ and let $A^{(i)}_{e,e'}\in\g$ denote the matrix tuple where the $i$-th matrix has a 1 at position $(e,e')$ and all other entries (also in all other matrices) are 0. Since these matrices form a basis of $\g$, it follows that
\[\g f=\linspan\{A^{(i)}_{e,e'}f\}.\]
For a tensor $f$ we define the support of $f$ as the set of monomials (i.e., standard basis tensors) for which $f$ has nonzero coefficient. For a linear subspace $V\subseteq\C^{w_1w_2}\otimes\C^{w_2w_3}\otimes\cdots\otimes\C^{w_dw_{d+1}}$ we define the support of $V$ as the union of the supports of all $f\in V$.
We write $e\cap e'=\varnothing$ to indicate that two edges $e$ and $e'$ do not share any vertex. We write $|e\cap e'|=1$ if they share exactly one vertex. We observe that for $f\in\{\fcom,f_0\}$ the vector space $T_f$ decomposes into a direct sum of three vector spaces,
\begin{align*}
\g_2&:=\linspan\{A^{(i)}_{e,e'}\mid1\le i\le d,1\le e,e'\le w_iw_{i+1},e\cap e'=\varnothing\},\\
\g_1&:=\linspan\{A^{(i)}_{e,e'}\mid1\le i\le d,1\le e,e'\le w_iw_{i+1},|e\cap e'|=1\},\\
\g_0&:=\linspan\{A^{(i)}_{e,e}\mid1\le i\le d,1\le e\le w_iw_{i+1}\},\\
\g&=\g_0\oplus\g_1\oplus\g_2,\\
T_f&=\g_0f\oplus\g_1f\oplus\g_2f.
\end{align*}
The last direct sum decomposition follows from the fact that $\g_0f$, $\g_1f$, and $\g_2f$ have pairwise disjoint supports.
We show in this section that $\dim\g_2\fcom=\dim\g_2f_0$, and that $\dim\g_1\fcom=\dim\g_1f_0$. In Section 8 we show that $\dim\g_0\fcom>\dim\g_0f_0$, which then implies $f_0\notin G\fcom$ by Claim 10. In fact, Theorem 13 gives the exact dimension of $\g_0\fcom$ by proving that $\g_0\fcom$ is isomorphic to the vector space of flows on the ABP digraph when identifying vertices in $V^1$ with their corresponding vertices in $V^{d+1}$. Theorem 14 establishes an additional equation based on the vertex parities that shows that $\g_0f_0$ is strictly lower dimensional than $\g_0\fcom$.
We start with Lemma 11, which shows that $\dim\g_2\fcom$ and $\dim\g_2f_0$ have full dimension.
\begin{lemma}
Let $f\in\{\fcom,f_0\}$. The space $\g_2f$ has full dimension. That is, its dimension equals $\sum_{i=1}^d w_iw_{i+1}(w_i-1)(w_{i+1}-1)$.
\end{lemma}
\begin{proof}
Suppose $f=\fcom$. The other case being analogous, we only argue this case.
We analyze the monomials that appear in the different $A^{(i)}_{e,e'}\fcom$ and argue that a monomial that appears in some $A^{(i)}_{e,e'}\fcom$ can only appear in that specific $A^{(i)}_{e,e'}\fcom$. Indeed, each monomial corresponds to a valid path in which one edge $e$ in layer $i$ is changed to $e'$. Since $e$ and $e'$ share no vertex, from this edge sequence we can reconstruct $i,e$, and $e'$ uniquely: $e'$ is the edge that does not have any vertex in common with the rest of the edge sequence, $i$ is its layer, and $e$ is the unique edge that we can replace $e'$ by in order to form a valid path. We conclude that the $A^{(i)}_{e,e'}\fcom$ have disjoint support and the lemma follows.
\end{proof}
To establish that $\dim\g_1\fcom=\dim\g_1f_0$, we introduce some notation.
For a connected directed graph $G=(V,E)$ we define a flow to be a labeling of the edge set $E$ by complex numbers such that at every vertex the sum of the labels of the incoming edges equals the sum of the labels of the outgoing edges. It is easily seen that the set of flows forms a vector space $F$. We have
\[\dim F=|E|-|V|+1,\tag{11}\label{eq:flowdim}\]
see e.g. Theorem 20.7 in [8].
Recall that $E^i$ denotes the set of edges from $V^i$ to $V^{i+1}$. Let $\mathcal X:=E^1\times\cdots\times E^d$ denote the direct product of the sets of edge lists. Each directed path of length $d$ from layer 1 to $d+1$ is an element of $\mathcal X$, but $\mathcal X$ contains other edge sets as well. Define $E_i:=\C^{E^i}$. Consider the following map $\varphi$ from $\mathcal X$ to $E_1\otimes\cdots\otimes E_d$,
\[\varphi(e_1,\ldots,e_d)=x_{e_1}\otimes\cdots\otimes x_{e_d}\in E_1\otimes\cdots\otimes E_d\]
where $(x_j)$ is the standard basis of $E_i$. Note $\varphi$ is a bijection between $\mathcal X$ and the standard basis of $E_1\otimes\cdots\otimes E_d$.
An edge set in $\mathcal X$ is called a valid path if it forms a path that starts and ends at corresponding vertices (see Sec.1). Let $\mathcal P\subseteq\mathcal X$ denote the set of valid paths.
\begin{proposition}
$\dim\g_1\fcom=\dim\g_1f_0=\sum_{i=1}^d(w_{i-1}+w_{i+1}-1)(w_i-1)w_i$, where $w_0:=w_d$.
\end{proposition}
\begin{proof}
The proof works almost analogously for $\fcom$ and $f_0$, so we treat only the more natural case $\fcom$. We show that $\g_1\fcom$ is isomorphic to a direct sum of vector spaces of flows on very simple digraphs. Fix $1\le i\le d$. Fix distinct $1\le a,b\le w_i$. For distinct edges $e,e'\in E^i$, let $\mathcal P_{e,e'}\subseteq\mathcal X$ be the set of edge sets containing $e'$ that are not valid paths, but that become valid paths by removing $e'$ and adding $e$. Let $\mathcal P^i_{a,b}\subseteq\mathcal X$ be the set of edge sets that are not valid paths, but that become valid paths by switching the end point of the $(i-1)$-th edge to $v_b^i$ and that also become valid paths by switching the start point of the $i$-th edge to $v_a^i$ (if $i-1=0$, then interpret $i-1:=d$). Pictorially, this means that elements in $\mathcal P^i_{a,b}$ are almost valid paths, but there is a discontinuity at layer $i$, where the path jumps from vertex $v_a^i$ to vertex $v_b^i$. We have
\[A^{(i)}_{e,e'}\fcom=\sum_{p\in\mathcal P_{e,e'}}\varphi(p).\]
The vectors $\{A^{(i)}_{e,e'}\fcom\mid1\le i\le d,e,e'\in E^i,|e\cap e'|=1\}$ are not linearly independent, because for $a\ne b$ we have
\begin{multline*}
\sum_{\substack{e,e'\text{ have same start point}\\e'\text{ ends at the }a\text{-th vertex}\\e\text{ ends at the }b\text{-th vertex}}}A^{(i-1)}_{e,e'}\fcom=\sum_{p\in\mathcal P^i_{a,b}}\varphi(p)\\
=\sum_{\substack{h,h'\text{ have same end point}\\h\text{ starts at the }a\text{-th vertex}\\h'\text{ starts at the }b\text{-th vertex}}}A^{(i)}_{h,h'}\fcom.\tag{12}\label{eq:dependency}
\end{multline*}
Define
\begin{align*}
T_{a,b,i}:={}&\linspan\left\{A^{(i-1)}_{e,e'}\fcom\;\middle|\;\begin{array}{l}e,e'\text{ have same start point}\\e'\text{ ends at the }a\text{-th vertex}\\e\text{ ends at the }b\text{-th vertex}\end{array}\right\}\\
&+\linspan\left\{A^{(i)}_{h,h'}\fcom\;\middle|\;\begin{array}{l}h,h'\text{ have same end point}\\h\text{ starts at the }a\text{-th vertex}\\h'\text{ starts at the }b\text{-th vertex}\end{array}\right\}.
\end{align*}
The support of $T_{a,b,i}$ and $T_{\widetilde a,\widetilde b,\widetilde i}$ are disjoint, provided $(a,b,i)\ne(\widetilde a,\widetilde b,\widetilde i)$. Hence
\[\g_1\fcom=\bigoplus_{\substack{1\le i\le d\\1\le a,b\le w_i\\a\ne b}}T_{a,b,i}.\]
It remains to prove that the dimension of $T_{a,b,i}$ is $w_{i-1}+w_{i+1}-1$, because then
\[\dim\g_1\fcom=\sum_{\substack{1\le i\le d\\1\le a,b\le w_i\\a\ne b}}(w_{i-1}+w_{i+1}-1)=\sum_{i=1}^d(w_{i-1}+w_{i+1}-1)(w_i-1)w_i.\]
Note that $T_{a,b,i}$ is defined as the linear span of $w_{i-1}+w_{i+1}$ many vectors, but (12) shows that these are not linearly independent. We prove that (12) is the only equality by showing that $T_{a,b,i}$ is isomorphic to a flow vector space. We define a multigraph with two vertices: $\odot$ and $\circledast$. We have $w_{i+1}$ many edges from $\odot$ to $\circledast$, and we have $w_{i-1}$ many edges from $\circledast$ to $\odot$. We denote by $\circledast\overset k\to\odot$ the $k$-th edge from $\circledast$ to $\odot$. Let $F_{a,b,i}$ denote the vector space of flows on this graph. Its dimension is $w_{i-1}+w_{i+1}-1$, see (11). We define $\varrho:E_1\otimes\cdots\otimes E_d\to F_{a,b,i}$ on rank 1 tensors via
\begin{align*}
\varrho(x_{e_1}\otimes\cdots\otimes x_{e_d})(\circledast\overset k\to\odot)&=\begin{cases}1&\text{if }e_{i-1}\text{ starts at }k\text{ in layer }i-1\\&\text{and ends at }a\text{ in layer }i,\\0&\text{otherwise},\end{cases}\\
\varrho(x_{e_1}\otimes\cdots\otimes x_{e_d})(\odot\overset l\to\circledast)&=\begin{cases}1&\text{if }e_i\text{ starts at }b\text{ in layer }i\\&\text{and ends at }l\text{ in layer }i+1,\\0&\text{otherwise}.\end{cases}
\end{align*}
Using (12) it is readily verified that $\varrho$ maps $T_{a,b,i}$ to $F_{a,b,i}$. It remains to show that $\varrho:T_{a,b,i}\to F_{a,b,i}$ is surjective. Let $\alpha:=|\mathcal P^i_{a,b}|$. We observe that
\begin{align*}
\varrho(A^{(i-1)}_{e,e'}\fcom)(\circledast\overset k\to\odot)&=\begin{cases}\alpha/w_{i-1}&\text{if }e,e'\text{ both start at the }k\text{-th vertex},\\0&\text{if }e,e'\text{ both start at the same vertex,}\\&\text{but not at the }k\text{-th},\end{cases}\\
\varrho(A^{(i-1)}_{e,e'}\fcom)(\odot\overset l\to\circledast)&=\alpha/(w_{i-1}w_{i+1}),\\
\varrho(A^{(i)}_{h,h'}\fcom)(\odot\overset l\to\circledast)&=\begin{cases}\alpha/w_{i+1}&\text{if }h,h'\text{ both end at the }l\text{-th vertex},\\0&\text{if }h,h'\text{ both end at the same vertex,}\\&\text{but not at the }l\text{-th},\end{cases}\\
\varrho(A^{(i)}_{h,h'}\fcom)(\circledast\overset k\to\odot)&=\alpha/(w_{i-1}w_{i+1}).
\end{align*}
Let $\Xi:=\sum A^{(i-1)}_{e,e'}\fcom$. Then $\forall k:\varrho(\Xi)(\circledast\overset k\to\odot)=\alpha/w_{i-1}$ and $\forall l:\varrho(\Xi)(\odot\overset l\to\circledast)=\alpha$. Therefore, for $e,e'$ starting at the $k_0$-th vertex and $h,h'$ ending at the $l_0$-th vertex we have that
\[\varrho\left(w_{i-1}w_{i+1}\varrho(A^{(i-1)}_{e,e'}\fcom)+w_{i-1}w_{i+1}\varrho(A^i_{h,h'}\fcom)-\Xi\right)\]
is nonzero only on exactly two edges: $\circledast\overset{k_0}\to\odot$ and $\odot\overset{l_0}\to\circledast$. Cycles form a generating set of the vector space $F_{a,b,i}$, which finishes the proof of the surjectivity of $\varrho$.
\end{proof}
